\chapter{Présentation du projet}

\section{Introduction}
Ce dossier présente le projet \textit{Badge reader}, un système de gestion de badges que j’ai conçu pour sécuriser et fluidifier les accès au bâtiment de Colint. En tant qu’étudiant en formation \textbf{Concepteur Développeur d’Applications (CDA)}, j’ai pu constater les limites du contrôle d’accès actuel, ce qui m’a motivé à développer cette solution.

Dans ce chapitre, je détaille mon rôle dans le développement d’un site web dédié à la gestion des badges, ainsi que l’organisation du projet et mes responsabilités. L’objectif est de répondre aux exigences du titre RNCP 37873, tout en proposant une solution adaptée aux besoins de l’école.
\newline

\begin{for_you}{Contexte et organisation du projet :}
    Dans le cadre de mon \textbf{projet de fin d’études (PFE)}, je mène la conception et le développement du système \textit{Badge reader}, de l’analyse des besoins à la réalisation finale. Ce projet s’adresse aux membres de Colint et, à terme, à d’autres interlocuteurs.
    \newline

    Le système repose sur l’attribution d’un \textbf{badge individuel} à chaque utilisateur, permettant de contrôler les entrées et sorties du bâtiment. Les accès sont définis par des \textbf{rôles} :
    \begin{itemize}
        \item Après 18h, les étudiants ne peuvent plus entrer (toute sortie est définitive).
        \item Le staff dispose d’un droit d’entrée jusqu’à 20h.
        \newline
    \end{itemize}
    
    Le déploiement est prévu pour le \textbf{1\textsuperscript{er} septembre 2026}.

    L’application sera développée en \textbf{Elixir} avec le framework \textbf{Phoenix}, pour une intégration cohérente avec l’intranet existant de l’école. L’organisation du projet s’appuie sur \textbf{GitHub Projects} et un \textbf{système Kanban} pour structurer les tâches et suivre leur progression.
\end{for_you}

\newpage

\section{Problématique et objectifs SMART}
\textbf{Problématique identifiée :}
Comment rendre l’accès au bâtiment de Colint \textbf{plus fluide pour les étudiants et le staff}, tout en garantissant sa \textbf{sécurité} et son \textbf{contrôle} ?
\newline

\textbf{Bénéfices attendus :}
\begin{itemize}
    \item Fluidifier le trafic dans le bâtiment.
    \item Vérifier la légitimité des personnes entrant dans l’école.
    \item Sécuriser les accès en fonction des rôles et des horaires.
    \newline
\end{itemize}

\begin{for_you}{Objectifs SMART et planning :}
    \textbf{\underline{Phase 1 – Application principale}} :
    \begin{itemize}
        \item Développer une application pour \textbf{collecter et exploiter 100\% des données} liées aux entrées/sorties.
        \item Gérer les utilisateurs, les rôles et les badges (physiques ou numériques).
        \item Livrer une \textbf{version 1.0 minimaliste et efficace} d’ici \textbf{mai 2026} (5 mois).
    \end{itemize}

    \textbf{\underline{Phase 2 – Badge dématérialisé}} :
    \begin{itemize}
        \item Remplacer les badges physiques par une \textbf{solution mobile} (iOS/Android).
        \item Intégrer cette fonctionnalité dans une \textbf{version 1.1} en \textbf{2 mois}.
        \newline
    \end{itemize}

    \textbf{\underline{Diagramme de contexte :}}
    \begin{center}
        \includegraphics[width=15cm,height=80mm]{./assets/uml/diagramme_de_context.png}
    \end{center}
\end{for_you}

\section{Liens utiles}
Les ressources ci-dessous m’ont servi de référence pour la méthodologie \textbf{SMART}, la gestion de projet et l’organisation du travail :
\newline

\begin{itemize}
    \item Documentation GitHub : \url{https://docs.github.com/}
    \item Objectifs SMART (Atlassian) : \url{https://bit.ly/smart-goals-atlassian}
    \item Project Management Institute : \url{https://www.pmi.org/}
    \item Manifeste Agile : \url{https://agilemanifesto.org/}
    \item Business Model Canvas : \url{https://bit.ly/business-model-canvas}
\end{itemize}